\exercisesection{}

\begin{enumerate}
\def\labelenumi{\arabic{enumi})}
\item
  Let E and F be Hilbert spaces and let u be a continuous linear map from E into F. Show that u is a Fredholm map if and only if the image of u is closed and the kernels of u and \(u^{*}\) are finite-dimensional.
\item
  Let \(p\) be a real number such that \(1 \leqslant p < +\infty\) . Let \(u\) be a diagonal endomorphism of \(\ell_{\mathrm{K}}^{p}(\mathbf{N})\) . Show that \(u\) is a Fredholm endomorphism if and only if it is a Riesz endomorphism. In particular, every diagonal Fredholm endomorphism of \(\ell_{\mathrm{K}}^{p}(\mathbf{N})\) has index zero.
\item
  \begin{enumerate}
  \def\labelenumii{\alph{enumii})}
  \tightlist
  \item
    There exists a Banach space E, a continuous automorphism u of E, and a Riesz endomorphism v of E such that \(u \circ v\) is not a Riesz endomorphism. There even exist such examples for which \(u \circ v - v \circ u\) has finite rank.
  \end{enumerate}
\end{enumerate}

\begin{enumerate}
\def\labelenumi{\alph{enumi})}
\setcounter{enumi}{1}
\tightlist
\item
  There exists a Banach space E, an automorphism u of E, and a finite-rank endomorphism h such that \(u + h\) is not a Riesz endomorphism.
\end{enumerate}

\begin{enumerate}
\def\labelenumi{\arabic{enumi})}
\setcounter{enumi}{3}
\tightlist
\item
  Let E be an infinite-dimensional complex Hilbert space.
\end{enumerate}

\begin{enumerate}
\def\labelenumi{\alph{enumi})}
\item
  The group \(\mathbf{U}(\mathrm{E})\) of unitary operators on E, equipped with the topology induced by the norm of \(\mathcal{L}(\mathrm{E})\) , is path-connected. (Show that if u is unitary, there exists a Hermitian operator v such that \(u = \exp(iv)\) .)
\item
  The group of linear isomorphisms of E is path-connected. (Use the polar decomposition, cf. I, p. 139, def. 3.)
\item
  For every \(n \in Z\) the set of Fredholm endomorphisms of E of index n is open and connected. It is nonempty.
\end{enumerate}

\begin{enumerate}
\def\labelenumi{\arabic{enumi})}
\setcounter{enumi}{4}
\tightlist
\item
  Let E and F be infinite-dimensional Banach spaces. A continuous linear map \(u: \mathrm{E} \to \mathrm{F}\) is said to be strictly singular if there does not exist a closed subspace \(\mathrm{F}_1 \subset \mathrm{E}\) of infinite dimension such that \(u\) induces, by passage to subspaces, an isomorphism from \(\mathrm{F}_1\) onto \(u(\mathrm{F}_1)\) .
\end{enumerate}

\begin{enumerate}
\def\labelenumi{\alph{enumi})}
\tightlist
\item
  Let u be a continuous linear map from E into F. The following conditions are equivalent:
\end{enumerate}

\begin{enumerate}
\def\labelenumi{(\roman{enumi})}
\item
  The map \(u\) is strictly singular;
\item
  For every closed subspace \(F_{1}\) of E of infinite dimension, there does not exist a real number \(\delta > 0\) such that \(\|u(x)\| \geqslant \delta\|x\|\) for all \(x \in F_{1}\) .
\end{enumerate}

\begin{enumerate}
\def\labelenumi{\alph{enumi})}
\setcounter{enumi}{1}
\item
  Let \(u \colon \mathrm{E} \to \mathrm{F}\) be a continuous linear map such that the image of \(u\) is closed and the kernel of \(u\) is finite-dimensional. If \(v \colon \mathrm{E} \to \mathrm{F}\) is strictly singular, then the image of \(u + v\) is closed and its kernel is finite-dimensional.
\item
  The set of strictly singular endomorphisms of E is a closed two-sided ideal of \(\mathcal{L}(\mathrm{E})\) .
\item
  Every compact linear map is strictly singular. If E and F are Hilbert spaces, then every strictly singular map \(E \rightarrow F\) is compact.
\item
  There exist Banach spaces E and F and a strictly singular non-compact map \(E \rightarrow F\) .
\end{enumerate}

\(f)\) Let \(p\) and \(r\) be real numbers such that \(1 \leqslant p < r\) . Every continuous linear map \(\ell_{\mathrm{K}}^{r}(\mathbf{N}) \to \ell_{\mathrm{K}}^{p}(\mathbf{N})\) is strictly singular. (Use exercise 10 of III, p. 106.)

\begin{enumerate}
\def\labelenumi{\arabic{enumi})}
\setcounter{enumi}{5}
\item
  Let E be a complex Hilbert space. Let u be a Fredholm endomorphism of E. If u is normal, then \(\text{ind}(u) = 0\) .
\item
  Let E and F be Banach spaces, and let u be an element of \(\mathcal{L}(\mathrm{E};\mathrm{F})\) . Prove that u is a Fredholm map if and only if the following conditions are satisfied:
\end{enumerate}

\begin{enumerate}
\def\labelenumi{(\roman{enumi})}
\tightlist
\item
  There exists a continuous seminorm p on E and \(C \geqslant 0\) such that the canonical map from E into the separated completion of \((E, p)\) is compact, and
\end{enumerate}

\[
\| x \| \leqslant \mathrm{C} \| u (x) \| + p (x)
\]

for all \(x \in E\) ;

\begin{enumerate}
\def\labelenumi{(\roman{enumi})}
\setcounter{enumi}{1}
\tightlist
\item
  There exists a continuous seminorm \(q\) on \(\mathrm{F}'\) and \(\mathrm{C}' \geqslant 0\) such that the canonical map from \(\mathrm{F}'\) into the separated completion of \((\mathrm{F}', q)\) is compact, and
\end{enumerate}

\[
\| y ^ {\prime} \| \leqslant \mathrm{C} ^ {\prime} \| ^ {t} u (y ^ {\prime}) \| + q (y ^ {\prime})
\]

for all \(y' \in F'\) .

\begin{enumerate}
\def\labelenumi{\arabic{enumi})}
\setcounter{enumi}{7}
\tightlist
\item
  Let U be a connected open subset of C. We denote by \(\mathcal{H}(\mathrm{U})\) the space of holomorphic functions on U equipped with the topology of compact convergence (EVT III, p. 10, example \(d\) ).
\end{enumerate}

\begin{enumerate}
\def\labelenumi{\alph{enumi})}
\item
  Let \(a \in \mathcal{H}(U)\) . Multiplication by a is a Fredholm endomorphism of \(\mathcal{H}(U)\) if and only if the set of zeros of a is finite. Then compute its index.
\item
  Let \(\mathcal{U} = (\mathrm{U}_{i})_{i \in \mathrm{I}}\) a covering of U by nonempty connected open sets. We set
\end{enumerate}

\[
\begin{array}{l} \mathrm{I} _ {2} = \{(i, j) \in \mathrm{I} \times \mathrm{I} | \mathrm{U} _ {i} \cap \mathrm{U} _ {j} \neq \varnothing \}, \\ \mathrm{I} _ {3} = \{(i, j, k) \in \mathrm{I} \times \mathrm{I} \times \mathrm{I} | \mathrm{U} _ {i} \cap \mathrm{U} _ {j} \cap \mathrm{U} _ {k} \neq \varnothing \}. \end{array}
\]

We consider the vector spaces

\[
\begin{array}{l} \mathrm{Z} ^ {1} (\mathcal {U}; \mathbf {C}) = \{(c _ {i, j}) \in \mathbf {C} ^ {\mathrm{I} _ {2}} | c _ {i j} + c _ {j k} + c _ {k i} = 0 \text { for all } (i, j, k) \in \mathrm{I} _ {3} \} \\ \mathrm{B} ^ {1} (\mathcal {U}; \mathbf {C}) = \{(c _ {i, j}) \in \mathbf {C} ^ {\mathrm{I} _ {2}} | \text { there exists } c ^ {\prime} \in \mathbf {C} ^ {\mathrm{I}} \text { such that } \end{array}
\]

\[
c _ {i j} = c _ {i} ^ {\prime} - c _ {j} ^ {\prime} \text {   for all   } (i, j) \in \mathrm{I} _ {2} \}.
\]

One has \(\mathrm{B}^1 (\mathcal{U};\mathbf{C})\subset \mathrm{Z}^1 (\mathcal{U};\mathbf{C})\) . We set

\[
\mathrm{H} ^ {1} (\mathcal {U}; \mathbf {C}) = \mathrm{Z} ^ {1} (\mathcal {U}; \mathbf {C}) / \mathrm{B} ^ {1} (\mathcal {U}; \mathbf {C}).
\]

Let \(\mathcal{V}=(\mathrm{V}_{\alpha})_{\alpha\in\mathrm{A}}\) be a covering of U by nonempty connected open sets, finer than U, and let \(\varrho\colon A\to I\) be a map such that \(V_{\alpha}\subset U_{\varrho(\alpha)}\) for all \(\alpha\in A\) .

The map \((c_{ij})_{i,j} \mapsto (c_{\varrho(\alpha)\varrho(\beta)})_{\alpha,\beta}\) induces a linear map, independent of the choice of \(\varrho\) , from \(\mathrm{H}^1(\mathcal{U};\mathbf{C})\) into \(\mathrm{H}^1(\mathcal{V};\mathbf{C})\) .

We denote by \(H^{1}(U; \mathbf{C})\) the inductive limit of vector spaces \(H^{1}(\mathcal{U}; \mathbf{C})\) equipped with these linear maps, as U ranges over the set of coverings of U by nonempty connected open sets.

\begin{enumerate}
\def\labelenumi{\alph{enumi})}
\setcounter{enumi}{2}
\tightlist
\item
  Let D be the endomorphism of \(\mathcal{H}(\mathrm{U})\) defined by \(Df = f'\) . Show that there exists an isomorphism of vector spaces from the cokernel of D onto \(\mathrm{H}^1 (\mathrm{U};\mathbf{C})\) . Deduce from this that D is a Fredholm endomorphism if and only if the space \(\mathrm{H}^1 (\mathrm{U};\mathbf{C})\) is finite-dimensional. Then compute the index of D.
\end{enumerate}

